\chapter{Introduction}
\label{chap:introduction}
\chapterrunning{introduction}

The design of effective heuristics remains one of the central practical problems in computer
science. Many computational tasks of scientific and industrial importance, including routing,
scheduling, and resource allocation, are too expensive to solve exactly at the scales demanded by
real applications. As a result, practitioners routinely rely on hand-crafted heuristics,
metaheuristics, or solver portfolios whose quality depends heavily on human insight, empirical
tuning, and domain knowledge \cite{rice1976,bengio2021mlco,burke2013,dokeroglu2024}. The
question of how to automate some portion of this design process has therefore persisted across
operations research, machine learning, and evolutionary computation for decades.

Recent progress in large language models (LLMs) changes the form of this question. Modern code
models can synthesize non-trivial executable programs from natural-language specifications, solve
benchmark programming tasks, and improve performance through repeated sampling and automated
testing \cite{chen2021codex,austin2021program,li2022alphacode,zan2023nl2code,jiang2024codesurvey,jain2024livecodebench}.
More importantly for optimization, these models can be embedded inside outer loops that evaluate
candidate programs, retain strong variants, and request improved descendants. Systems such as
AlphaDev and FunSearch show that evaluator-guided search can go beyond naive one-shot completion
when the problem admits fast verification and a suitably structured search space
\cite{mankowitz2023alphadev,romeraparedes2024funsearch}. At the same time, early work on LLMs as
evolutionary optimizers suggests both promise and sharp scale limitations, especially when the
search is applied to combinatorial optimization rather than small synthetic benchmarks
\cite{liu2023lmea,daros2025llmco}.

These developments motivate the present dissertation. The central question is whether repeated
proposal, evaluation, and revision can improve executable heuristics in a way that survives
stronger benchmarks and stronger baselines. That question also requires tighter controls on
interpretation. Apparent gains may
reflect useful adaptation, but they may also reflect benchmark headroom, weak baselines, or prompt
sequences that generate many variants without preserving a real improvement. The dissertation
examines iterative code evolution as an empirical object that must be measured across more than one
family and judged by held-out results rather than by isolated successful artifacts.

\section{Problem Statement}
\label{sec:intro-problem}

This dissertation studies LLM-guided code evolution as a \emph{measurable adaptive search loop}
rather than as a general claim of autonomous algorithm invention. The framework considered here is
simple by design. A model receives a constrained programming interface and task feedback, proposes
a deterministic Python heuristic, and then receives structured information about the heuristic's
behaviour and performance. The loop repeats over many epochs, storing code variants, validator
signals, and archive summaries that can be replayed into later prompts. The scientific objective
is to determine when such a loop adapts in a useful way and when it merely produces syntactically
different but functionally unimportant variants.

The dissertation addresses the following research question:

\begin{quote}
Under what conditions does an LLM-guided code-evolution loop generate useful heuristic adaptation,
and when are apparent gains better explained by search budget, benchmark headroom, validator
pressure, or baseline strength?
\end{quote}

This broad question is instantiated through several narrower empirical questions distributed
across the dissertation:

\begin{enumerate}[leftmargin=2em]
\item Do replay-aware memory mechanisms improve held-out performance relative to no replay, random
replay, and other replay-selection rules?
\item Does the answer remain stable when the environment shifts from toy adversarial games to
standard benchmark optimization families such as TSPLIB and CVRPLIB?
\item Can the search loop discover reusable heuristic components or adaptive controllers, rather
than only full-program variants tied to a single scaffold?
\item How strongly do validator pressure, baseline strength, and benchmark headroom limit the gains
that the loop can realize on held-out instances?
\end{enumerate}

\section{Research Objectives}
\label{sec:intro-objectives}

The project is organized around four concrete objectives.

\begin{enumerate}[leftmargin=2em]
\item To formalize LLM-guided code evolution as a controlled experimental pipeline rather than an
anecdotal prompt-engineering exercise.
\item To evaluate the pipeline on benchmark families whose structure and baseline difficulty differ
materially, namely simple games, TSP, ATSP, and CVRP.
\item To distinguish family-specific replay effects from more general claims about memory, broader
search, or benchmark-specific headroom.
\item To produce a conservative account of the present capabilities and limits of executable
heuristic evolution by LLMs, including null results and bounded wins.
\end{enumerate}

\section{Working Hypotheses}
\label{sec:intro-hypotheses}

The hypotheses below state the principal empirical claims examined in this dissertation. They are
framed to distinguish mechanism-specific effects from broader effects due to search budget,
benchmark structure, and baseline strength.

\begin{enumerate}[leftmargin=2em]
\item \textbf{H1: Adaptive search under validator feedback.} When a benchmark family contains
reachable heuristic headroom and the interface remains executable under repeated sampling, an
iterative LLM-guided loop will generate non-trivial behavioural diversity and occasionally improve
held-out quality over one-shot sampling.
\item \textbf{H2: Replay must justify its added machinery.} A replay or archive mechanism should be
credited only when it improves held-out performance relative to simpler controls or alternative
replay rules evaluated under the same protocol.
\item \textbf{H3: Gains shrink under stronger evaluators.} Apparent advantages of evolved code
should weaken as validator pressure rises, feasibility constraints tighten, and classical baselines
consume more of the available benchmark headroom.
\item \textbf{H4: Transfer is partial rather than universal.} Search components discovered on one
family may transfer across related problems, but the strength of that transfer should depend on the
interaction structure of the target domain rather than on a task-agnostic notion of ``general
algorithmic intelligence.''
\end{enumerate}

These hypotheses permit negative or mixed outcomes. Under controlled evaluation, null results are
informative because they help separate genuine mechanism effects from artefacts of sampling budget
or benchmark choice.

\section{Formal Setup and Notation}
\label{sec:intro-formal-setup}

The dissertation treats heuristic search as a problem of learning or discovering \emph{programs}
that map benchmark instances to feasible solutions. Let $\mathcal{X}$ denote a family of problem
instances, let $\Omega(x)$ denote the feasible solution set for instance $x \in \mathcal{X}$, and
let $f(\cdot;x)$ denote a minimization objective on $\Omega(x)$.

\begin{definition}
A \emph{benchmark optimization family} is a triple
$\mathcal{F} = (\mathcal{X}, \Omega, f)$ consisting of an instance distribution or benchmark set
$\mathcal{X}$, a feasibility map $\Omega$, and an objective function $f$ to be minimized.
\end{definition}

\begin{definition}
A \emph{deterministic heuristic program} for $\mathcal{F}$ is an executable mapping
$h : \mathcal{X} \to \Omega(x) \cup \{\bot\}$ that returns either a feasible solution or a
detectable failure symbol $\bot$ under a fixed interface and resource budget.
\end{definition}

This definition covers the constructive game agents of the early phases, the TSP and ATSP
tour-construction heuristics of the routing phases, and the bounded CVRP solvers evolved in the
later real-world experiments. The relevant distinction is whether a program yields a measurable
objective value under controlled execution.

For any evaluation panel $\mathcal{D} \subseteq \mathcal{X}$, the empirical quality of a heuristic
program is summarized by an aggregate metric
\begin{equation}
J_{\mathcal{D}}(h) = \frac{1}{|\mathcal{D}|}\sum_{x \in \mathcal{D}} \operatorname{gap}(h,x),
\label{eq:intro-aggregate-gap}
\end{equation}
where lower values are better and the instance-level gap is defined, whenever a reference cost
$f^{\star}(x)$ is available, by
\begin{equation}
\operatorname{gap}(h,x) =
\frac{f(h(x);x)-f^{\star}(x)}{f^{\star}(x)}.
\label{eq:intro-relative-gap}
\end{equation}
For feasibility-constrained settings such as CVRP, the implementation can replace $f$ by a
penalized objective $\tilde{f}$ that incorporates route-capacity or constraint violations. The
form of \eqref{eq:intro-relative-gap} is standard in routing benchmarks because it allows
comparisons across instances of different scales.

\begin{definition}
An \emph{LLM-guided code-evolution loop} is an iterative process over a program space
$\mathcal{H}$ with prompt state $p_t$, archive state $\mathcal{A}_t$, and evaluator $E$, such that
\begin{align}
h_t &\sim q_{\phi}(\cdot \mid p_t), \label{eq:intro-sampling}\\
r_t &= E(h_t; \mathcal{D}_{\mathrm{train}}), \label{eq:intro-evaluation}\\
\mathcal{A}_t &= U(\mathcal{A}_{t-1}, h_t, r_t), \label{eq:intro-archive-update}\\
p_{t+1} &= C(\mathcal{T}, \mathcal{A}_t, r_t), \label{eq:intro-prompt-update}
\end{align}
where $q_{\phi}$ is the model-conditioned program generator, $U$ is the archive update rule, and
$C$ constructs the next prompt from the task specification $\mathcal{T}$ and archived evidence.
\end{definition}

Equations \eqref{eq:intro-sampling}--\eqref{eq:intro-prompt-update} make the research claim
precise: replay is not a vague intuition, but a specific modification of the archive state
$\mathcal{A}_t$ and prompt-construction operator $C$. A replay mechanism has independent value only
if it improves final held-out performance relative to an explicit control condition. One generic
comparison can be expressed as
\begin{equation}
\Delta_{\mathrm{replay}} =
J_{\mathcal{D}_{\mathrm{test}}}\!\left(h^{\mathrm{replay}}\right) -
J_{\mathcal{D}_{\mathrm{test}}}\!\left(h^{\mathrm{control}}\right),
\label{eq:intro-replay-delta}
\end{equation}
where a robustly negative value of $\Delta_{\mathrm{replay}}$ would support a replay-specific
advantage on a lower-is-better metric.

To separate meaningful behavioural diversity from mere lexical change, the loop also tracks
novelty-like signals. A generic archive-relative novelty score can be written as
\begin{equation}
\nu(h_t,\mathcal{A}_{t-1}) =
1 - \max_{h \in \mathcal{A}_{t-1}} \operatorname{sim}(h_t,h),
\label{eq:intro-novelty}
\end{equation}
where $\operatorname{sim}$ is a code- or behaviour-level similarity measure. No single novelty
metric is assumed to be canonical; novelty is treated as a diagnostic that must be interpreted
jointly with validator outcomes and held-out quality.

\begin{table}[t]
\centering
\small
\setlength{\tabcolsep}{3pt}
\begin{tabular}{>{\raggedright\arraybackslash}p{0.95in} >{\raggedright\arraybackslash}p{3.85in}}
\toprule
Symbol & Meaning \\
\midrule
$\mathcal{X}$ & benchmark family or instance space \\
$x \in \mathcal{X}$ & a specific optimization instance \\
$\Omega(x)$ & feasible solution set for instance $x$ \\
$f(\cdot;x)$ & instance-level minimization objective \\
$J_{\mathcal{D}}(h)$ & aggregate evaluation metric of heuristic $h$ on panel $\mathcal{D}$ \\
$\mathcal{H}$ & constrained program space searched by the outer loop \\
$q_{\phi}$ & model-conditioned proposal distribution over candidate programs \\
$\mathcal{A}_t$ & archive or replay state available at iteration $t$ \\
$E$ & validator and scorer used to evaluate a candidate program \\
$\nu(h_t,\mathcal{A}_{t-1})$ & archive-relative novelty diagnostic for candidate $h_t$ \\
\bottomrule
\end{tabular}
\caption[Core notation]{Core notation used throughout the introduction and later methodological chapters.}
\label{tab:intro-notation}
\end{table}

\begin{figure}[t]
\centering
\begin{tikzpicture}[
  node distance=0.75cm and 0.6cm,
  box/.style={draw, rounded corners, align=center, minimum width=3.15cm, minimum height=0.95cm},
  line/.style={-Latex, thick}
]
\node[box] (task) {Task specification\\and interface constraints};
\node[box, below=of task] (prompt) {Prompt constructor\\$C(\mathcal{T}, \mathcal{A}_t, r_t)$};
\node[box, below=of prompt] (model) {LLM proposal step\\$h_t \sim q_{\phi}(\cdot \mid p_t)$};
\node[box, below left=of model, xshift=-0.1cm] (eval) {Validator and scorer\\$r_t = E(h_t;\mathcal{D}_{\mathrm{train}})$};
\node[box, below right=of model, xshift=0.1cm] (archive) {Archive update\\$\mathcal{A}_{t+1}=U(\mathcal{A}_t,h_t,r_t)$};
\draw[line] (task) -- (prompt);
\draw[line] (prompt) -- (model);
\draw[line] (model) -- (eval);
\draw[line] (model) -- (archive);
\draw[line] (eval) -- (archive);
\draw[line] (archive.north) |- (prompt.east);
\end{tikzpicture}
\caption[Conceptual code-evolution loop]{Conceptual structure of the LLM-guided code-evolution loop studied in this dissertation.
The central question is whether the closed loop between archive, validator, and proposal model
produces reliable adaptation under controlled protocols.}
\label{fig:intro-loop}
\end{figure}

\section{Why This Question Matters}
\label{sec:intro-motivation}

The practical motivation follows directly from the role of heuristics in large-scale optimization.
If language models can participate productively in heuristic search, then they could become useful
assistants for constructing bounded optimization procedures in domains where exact methods are too
slow and hand-designed heuristics are expensive to develop. The scientific motivation is subtler.
Recent work on LLM-based scientific and algorithmic discovery has highlighted competitive
programming performance, small algorithmic improvements, and evaluator-guided discoveries in
specialized settings
\cite{li2022alphacode,mankowitz2023alphadev,romeraparedes2024funsearch}. Those results are
important, but they do not by themselves establish a general understanding of when an LLM-driven
search loop should work, when it should fail, and which controls are necessary to distinguish true
adaptation from a larger sampling budget.

This issue is especially important in stochastic learning systems, where optimistic conclusions can
be driven by best-run anecdotes, fragile benchmarks, or insufficient controls. Empirical machine
learning has already seen the cost of such under-specified reporting, particularly in reinforcement
learning, where replicated evaluation, uncertainty-aware summaries, and strong baselines are now
widely recognized as necessary for credible claims \cite{henderson2018drltm,agarwal2021precipice}.
The same discipline is needed here. A dissertation on LLM-guided heuristic evolution should not
only report positive cases. It should also identify limits, null results, and failure mechanisms.

\section{Dissertation Scope}
\label{sec:intro-scope}

The work reported here begins with simple adversarial games and then moves to increasingly
realistic combinatorial optimization families. This progression is deliberate. The initial game
environments provide a cheap setting in which iterative adaptation, curriculum effects, replay,
and behavioural diversity can be observed directly. Later phases retain the same basic
code-evolution machinery but replace the toy environment with benchmark-driven tasks whose solution
quality can be evaluated against established references and classical heuristics. In particular,
the dissertation studies:

\begin{enumerate}[leftmargin=2em]
\item adversarial code evolution in simple two-player games;
\item replay-aware heuristic evolution on symmetric TSPLIB traveling salesperson benchmarks;
\item transfer of the same design ideas to asymmetric TSPLIB ATSP and CVRPLIB CVRP families;
\item modular operator discovery and adaptive heuristic portfolio control on TSP;
\item bounded whole-solver evolution on real-world CVRPLIB X instances; and
\item a cross-family synthesis of the completed simple-game, TSP, and real-world CVRP evidence,
together with a descriptive direct Codex baseline on bounded CVRP.
\end{enumerate}

\begin{table}[t]
\centering
\scriptsize
\setlength{\tabcolsep}{3pt}
\begin{tabular}{
  >{\raggedright\arraybackslash}p{0.65in}
  >{\raggedright\arraybackslash}p{1.1in}
  >{\raggedright\arraybackslash}p{1.35in}
  >{\raggedright\arraybackslash}p{1.1in}
}
\toprule
Family & Benchmark basis & Evolved artifact & Primary held-out metric \\
\midrule
Simple games & synthetic adversarial environments & deterministic game-playing policy code & win rate / score margin \\
TSP & TSPLIB symmetric routing instances & tour-construction or improvement heuristic code & relative optimality gap \\
ATSP & TSPLIB asymmetric routing instances & transfer-tested heuristic code & relative optimality gap \\
CVRP & CVRPLIB, especially Uchoa X instances & whole solver or route-construction code & penalized gap with feasibility pressure \\
\bottomrule
\end{tabular}
\caption[Problem families and evaluation focus]{Problem families covered by the dissertation and the corresponding evaluation focus.}
\label{tab:intro-families}
\end{table}

The study therefore extends beyond a single TSP setting, but it does not attempt to formulate a
universal theory of automatic algorithm discovery. Its aim is narrower: to develop a unified
framework for studying \emph{capabilities and limits} in LLM-guided code evolution across domains
with different structural demands.

\section{Main Thesis}
\label{sec:intro-main-thesis}

The central thesis of this dissertation is as follows:

\begin{quote}
Across simple games, TSP, ATSP, and real-world CVRP, LLM-guided code evolution can generate
diverse executable heuristic programs and sometimes preserve useful adaptations, but its success is
capability-bounded. In the current evidence, performance depends strongly on validator pressure,
baseline strength, benchmark headroom, and task-specific interaction structure. Replay-specific
mechanisms matter in some settings, but no single replay recipe remains dominant across the
completed benchmark families.
\end{quote}

The completed results support a conservative interpretation. They do not support two stronger
claims: that replay uniformly improves performance, or that iterative LLM search has already
surpassed mature human-designed optimization heuristics in a general sense. Instead, the empirical
record identifies where adaptation occurs, where it fails to transfer, and how meaningful
improvement can be separated from superficial novelty and favorable benchmark conditions.

\section{Research Contributions}
\label{sec:intro-contributions}

The dissertation makes five main contributions.

\begin{enumerate}[leftmargin=2em]
\item It introduces a unified experimental framework for iterative code evolution with archived
feedback, constrained execution, deterministic validators, and report-first research artifacts.
\item It extends the study of replay-aware code evolution from toy adversarial settings to multiple
benchmark optimization families, including symmetric TSP, asymmetric TSP, and capacitated vehicle
routing.
\item It evaluates not only positive search variants but also mechanism-focused alternatives,
including curated failure replay, compression pressure, modular operator evolution, and
portfolio-based selection under strong fixed baselines.
\item It emphasizes replicated, uncertainty-aware, benchmark-based evaluation in the spirit of
careful empirical RL methodology \cite{henderson2018drltm,agarwal2021precipice}.
\item It contributes a conservative cross-family synthesis in which negative and boundary results
are treated as findings, not as failures to be hidden.
\end{enumerate}

\section{Preview of Findings}
\label{sec:intro-preview}

The main empirical findings are mixed but coherent.

First, the game phases show real functional adaptation, but the novelty audit also shows that many
large code changes are mixed or negative rather than uniformly useful. Second, replay can help on
benchmark routing, but not in the simple way initially hypothesized. On symmetric TSP, naive
raw-failure replay is not the strongest mechanism; curation matters more than failure accumulation
alone, and compression does not help the strongest replay arms. Third, the family extensions block
any universal replay claim: ATSP yields only a weak compression-aware signal, while no-replay
remains strongest on the replicated phase-5 CVRP suite. Fourth, the later higher-upside phases act
as strict falsification tests. Modular operator discovery produces genuine rediscovery structure,
yet no operator survives the full validation pipeline. Adaptive portfolio control improves over a
weak one-shot LLM selector, yet does not beat the strongest fixed heuristic or the supervised
selector when oracle headroom is absent on the primary endpoint. Whole-solver evolution on
real-world CVRP stays feasible and improves over weaker bounded baselines, but not over the fixed
Clarke--Wright baseline. Finally, the cross-family synthesis supports a capability-bounded account
in which executable variation is common, but durable gains depend on the target family and on the
strength of the competing baseline.

\begin{figure}[t]
\centering
\begin{tikzpicture}
\begin{groupplot}[
  group style={group size=2 by 1, horizontal sep=1.1cm},
  width=0.40\textwidth,
  height=0.24\textheight,
  ymin=0,
  scaled y ticks=false,
  enlarge x limits=0.18,
  axis line style={black},
  tick style={black},
  ylabel={Lower-is-better gap},
  xtick=data,
  x tick label style={font=\scriptsize, rotate=25, anchor=east},
  yticklabel style={
    font=\footnotesize,
    /pgf/number format/fixed,
    /pgf/number format/precision=2
  },
  every axis title/.style={font=\small},
  title style={at={(0.5,1.03)}, anchor=south},
  point meta=y,
  nodes near coords={\pgfmathprintnumber[fixed,precision=3]{\pgfplotspointmeta}},
  every node near coord/.append style={font=\tiny, anchor=south, yshift=1pt, text=black},
]
\nextgroupplot[
  title={TSP replay suite},
  ymax=0.145,
  symbolic x coords={No replay,Failure,Compression,Random},
]
\addplot+[ybar, draw=black, fill=black!20, bar width=9pt] coordinates {
  (No replay,0.132474)
  (Failure,0.130918)
  (Compression,0.124895)
  (Random,0.104243)
};

\nextgroupplot[
  title={CVRP held-out suite},
  ymax=0.52,
  ylabel={},
  symbolic x coords={CW savings,Evolution,NN,Regret+LS},
]
\addplot+[ybar, draw=black, fill=black!45, bar width=9pt] coordinates {
  (CW savings,0.073766)
  (Evolution,0.182231)
  (NN,0.273435)
  (Regret+LS,0.466391)
};
\end{groupplot}
\end{tikzpicture}
\caption[Representative held-out phase results]{Representative held-out results from completed project phases. Lower is better in both
panels. In the TSP replay suite, random replay achieved the best mean transfer gap among the
replay variants. In the real-world CVRP suite, solver evolution improved over weaker bounded
baselines yet remained behind the fixed Clarke--Wright savings heuristic on the primary endpoint.}
\label{fig:intro-preview-results}
\end{figure}

Overall, these findings support a \emph{capability-and-limits} interpretation. The system is good
enough to warrant serious study, but not strong enough to justify claims of general
algorithm-discovery superiority.

\section{Organization of the Dissertation}
\label{sec:intro-organization}

The remainder of the dissertation is organized as follows. Chapter~\ref{chap:literature-review}
reviews the literature on heuristics and algorithm selection, machine learning for combinatorial
optimization, LLM-based code generation, evaluator-guided algorithm discovery, and curriculum,
replay, and diversity mechanisms. Chapter~\ref{chap:methodology} presents the common experimental
framework, replication discipline, and reporting methodology. Chapter~\ref{chap:simple-games}
reports the adversarial transfer evidence and the manual novelty adjudication. Chapter~
\ref{chap:routing-studies} examines the replicated routing replay and replay-mechanism results
across TSP, ATSP, and CVRP. Chapter~\ref{chap:advanced-phases} presents the operator-discovery,
adaptive-portfolio, and real-world CVRP solver-evolution phases. Chapter~\ref{chap:conclusion}
returns to the main hypotheses, synthesizes the cross-family evidence, and outlines the principal
limitations and future research directions.
